%\iffalse meta-comment
% Copyright (C) 2013 by Balogh Gergő
% This file and its companion matrials are free to use for all people associated with SZESZT (Szegedi Egyetemistákat Szocializáló Társaság). All other rights are reserved.

% This product is provided 'as is', means use at your on risk.
%\fi

%\iffalse
%<article|short|package>\NeedsTeXFormat{LaTeX2e}
%<article>\ProvidesClass{szesztarticle}
%<short>\ProvidesClass{szesztshort}
%<package>\ProvidesPackage{szeszt}
%<article|short|package>[2013/18/05 v1.0tinco szeszt]
%<article>\LoadClass{article}
%<short>\LoadClass[11pt]{article}
%
%<*driver>
\documentclass{ltxdoc}
\usepackage{szeszt}
\EnableCrossrefs
\CodelineIndex
\RecordChanges
\begin{document}
\DocInput{szeszt.dtx}
\end{document}
%</driver>
%\fi
%
%\changes{v0.0tinco}{2013/18/05}{Initial version}
%
%\GetFileInfo{szeszt.sty}
%
% \title{The \textsf{szeszt} package and class boundle\thanks{This document
% corresponds to \textsf{szeszt}~\fileversion,
% dated~\filedate.}}
% \author{Balogh Gerg\H{o}\\\texttt{geryxyz@gmail.com}}
%
% \maketitle
%
%    \begin{macrocode}
%<*short|article>
\RequirePackage{fancyhdr}
%</short|article>
%<*short>
\fancypagestyle{plain}{
    \fancyhead[L]{}
    \fancyhead[C]{}
    \fancyhead[R]{}
    \fancyfoot[L]{}
    \fancyfoot[C]{}
    \fancyfoot[R]{}
    \renewcommand{\headrulewidth}{0pt}
    \renewcommand{\footrulewidth}{0pt}
}
\fancyhead[L]{}
\fancyhead[C]{}
\fancyhead[R]{}
\fancyfoot[L]{}
\fancyfoot[C]{}
\fancyfoot[R]{}
\renewcommand{\headrulewidth}{0pt}
\renewcommand{\footrulewidth}{0pt}
\pagestyle{fancy}

\RequirePackage{calligra}
\RequirePackage{cantarell}
\RequirePackage{caption}
\captionsetup[figure]{labelformat=empty}
\RequirePackage{booktabs}
%</short>

\RequirePackage{calc}
\RequirePackage{geometry}
\RequirePackage{graphicx}
\RequirePackage{ifpdf}
\ifpdf %if using pdfLaTeX in PDF mode
    \DeclareGraphicsExtensions{.pdf,.png,.jpg,.jpeg,.mps}
\else %if using LaTeX or pdfLaTeX in DVI mode
    \DeclareGraphicsExtensions{.eps,.bmp}
    \DeclareGraphicsRule{.emf}{bmp}{}{}% declare EMF filename extension
    \DeclareGraphicsRule{.png}{bmp}{}{}% declare PNG filename extension
\fi
\RequirePackage{xcolor}
\RequirePackage{array}
\newsavebox{\titleright}
%    \end{macrocode}
%
%\section{Page level format}
%\begin{macro}{\bemargingood}
%Sets the margin for two side printing. |\bemargingood|\marg{unit}\marg{binding offset}
%    \begin{macrocode}
%<*article|package>
\newcommand{\bemargingood}[2]{
    \geometry{twoside,inner=#1*2,top=#1*3,outer=#1*4,bottom=#1*5,bindingoffset=#2}
}
\bemargingood{1cm}{20mm}
%</article|package>
%<*short>
\geometry{hmargin=2cm,top=1cm,nofoot,nohead}
\setlength{\leftskip}{0pt plus.5fil}
\setlength{\rightskip}{0pt plus-.5fil}
\setlength{\parfillskip}{0pt plus1fil}
\setlength{\parindent}{0pt}
\setlength{\parskip}{1ex}
%</short>
%    \end{macrocode}
%\end{macro}
%
%\section{Text level format}
%\begin{macro}{\maketitle}
%\begin{macro}{\szesztmaketitle}
%Simple title. Example:\\
%\title{Title of Example}
%\author{Balogh Gerg\H{o}}
%\date{2013. may 18.}
%\szesztmaketitle
%    \begin{macrocode}
%<*article>
\renewcommand{\maketitle}{
%</article>
%<*package>
\newcommand{\szesztmaketitle}{
%</package>
%<*article|package>
    \savebox{\titleright}{
        \parbox[b]{.5\textwidth}{
            \raggedright
            \textcolor{black!50!white}{\textit{\@author}}\vspace{1ex}\\
            \textbf{\textsf{\Large\@title}}\vspace{1.5ex}\\
            \textcolor{black!50!white}{\textit{\@date}}\\
        }
    }
    \noindent
    \parbox{\textwidth}{
        \includegraphics[height=.3\textwidth]{res/pic/logo}
        \rule{1pt}{.3\textwidth}
        \raisebox{.15\textwidth-.5\ht\titleright}{\usebox{\titleright}}
    }
    \vspace*{.15\textwidth}
}
%</article|package>
%<*short>
\renewcommand{\maketitle}{
    \begin{center}
        \textit{\@author}\vspace{1ex}\\
        {\hspace*{-2.5ex}\Huge\calligra\textcolor{black!25!white}{\@title}}\vspace{-3.5ex}\\
        \textbf{\textsf{\LARGE\@title}}\vspace{1.5ex}\\
        \textit{\@date}\vspace*{4ex}\\
    \end{center}
}
\newcommand{\maketitlenodate}{
    \begin{center}
        \textit{\@author}\vspace{1ex}\\
        {\hspace*{-2.5ex}\Huge\calligra\textcolor{black!25!white}{\@title}}\vspace{-3.5ex}\\
        \textbf{\textsf{\LARGE\@title}}\vspace*{4ex}\\
    \end{center}
}
%</short>
%    \end{macrocode} 
%\end{macro}
%\end{macro}
%\begin{macro}{\szesztdividingsection}
%Use to create a section title which separate two distinc part of the same document. For example think of the list of founders.
%    \begin{macrocode}
%<*short>
\newcommand{\szesztdividingsection}[1]{
    \begin{center}
        {\hspace*{-2.5ex}\huge\calligra\textcolor{black!25!white}{#1}}\vspace{-2.5ex}\\
        \textbf{\textsf{\large{}#1}}\vspace{8ex}\\
    \end{center}
}
%</short>
%    \end{macrocode}
%\end{macro}
%\begin{macro}{\szesztsignum}
%    \begin{macrocode}
\newcommand{\szesztsignum}[2][]{
    
    \parbox{.3\textwidth}{
        \begin{center}
            \textcolor{black!25!white}{\rule{.3\textwidth}{.1pt}}\\
            \textsf{#2}\vspace{1ex}\\
            {\scriptsize\textcolor{black!50!white}{
                    \begin{tabular}{>{\itshape\rule{0pt}{2.25ex}}c}
                        #1
                    \end{tabular}
                }
            }
        \end{center}
        \vspace*{2ex}
    }
}
%    \end{macrocode}
%\end{macro}
%\begin{macro}{szesztsignums}
%    \begin{macrocode}
\newenvironment{szesztsignums}{
    \begin{center}
    \begin{tabular}{@{}c@{\hspace*{.15\textwidth}}c@{}}    
}{
    \end{tabular}
    \end{center}
}
%    \end{macrocode}
%\end{macro}
%\begin{macro}{\szeszttableheader}
%    \begin{macrocode}
\newcommand{\szeszttableheader}[1]{
    \textbf{#1}
}
%    \end{macrocode}
%\end{macro}
%\begin{macro}{\szeszthighemph}
%    \begin{macrocode}
\newcommand{\szeszthighemph}[1]{
    \begin{center}
    \textbf{\fcafamily #1}
    \end{center}
}
%    \end{macrocode}
%\end{macro}
%\begin{macro}{\szesztdate}
%    \begin{macrocode}
\newcommand{\szesztdate}[1]{
    \begin{flushright}
        \textit{#1, \@date}
    \end{flushright}
}
%    \end{macrocode}
%\end{macro}
%\begin{macro}{\szeszt}
%    \begin{macrocode}
\newcommand{\szeszt}{
    SZESZT
}
%    \end{macrocode}
%\end{macro}
% \Finale
\endinput
